\chapter*{Executive Summary}
\addcontentsline{toc}{chapter}{Executive Summary}

\section*{Project Motivation and Goal}
AI-assisted development can introduce changes across dependency manifests, application code, and infrastructure-as-code. A change may be syntactically valid while weakening a security control, changing the logical resource that a policy protects, or leaving insufficient evidence to make a trustworthy decision. ASENT (Agent Sentinel) is being developed as an assurance gateway between a proposed AI-assisted change and the point where that change is trusted, merged, or deployed. Its goal is to aggregate evidence from relevant security checks into a deterministic decision of \texttt{ACCEPT}, \texttt{REVIEW}, or \texttt{BLOCK}.

\section*{System Scope}
ASENT is organized around three security boundaries: CAVR addresses dependency security, SATRA-RV addresses application-code security, and SABLE addresses infrastructure security expressed through Terraform. The shared workflow combines module evidence and applies a final gate. This report's strongest verified implementation and contribution evidence concerns SABLE and the shared evidence/final-gate path. It does not treat the presence of other modules in the repository as proof of their authorship or complete readiness.

\section*{Implemented SABLE Approach}
SABLE analyzes supplied Terraform/\allowbreak{}HCL snapshots locally, builds a bounded resource/\allowbreak{}reference model, checks a registered baseline obligation, identifies a supported successor for the protected logical asset, and evaluates a bounded IAM/\allowbreak{}S3 authorization model against that successor. The outcome is explicit: \texttt{PRESERVED} when the registered obligation continues to hold, \texttt{REGRESSED} when supported evidence shows that it does not hold, or \texttt{UNKNOWN} when input semantics or resource correspondence are insufficient. In the shared workflow, current evidence is associated with the candidate snapshot and checked for freshness, input-\allowbreak{}hash consistency, and integrity before it contributes to the final decision.

\section*{Current Evidence}
PR~\#5 integrated the CAVR, SATRA-RV, and SABLE work into \texttt{master} at merge commit \texttt{ffe8aa91}; its exact tested head was \texttt{554b6d51}, and all 18 checks passed on that head. Representative Core, SATRA, and SABLE workflow runs are linked in Chapter~7. Separately, the focused SABLE/shared-assurance CI run on source commit \texttt{84f91a9} passed 48 tests, and the local focused selection on the correction-containing checkout passed 27 tests with 10 deselected. These results belong to their identified commits and are not a full-suite pass.

The 2026-10-12 post-merge shared-dashboard demonstration recorded SABLE \texttt{PRESERVED} in run \texttt{21adc601d6b7}, with overall \texttt{REVIEW} because SATRA was \texttt{INCONCLUSIVE}; the wrong-target run \texttt{0849603bfc16} recorded SABLE \texttt{REGRESSED} and overall \texttt{BLOCK}. Both traces showed valid integrity and matching candidate/input bindings. Their original screenshots are preserved under \path{report/evidence}; no AWS resources were applied or queried. This is a two-scenario demonstration, not full acceptance of every module.

An earlier rehearsal on \texttt{sable} commit \texttt{753a55b} recorded preserved, regressed, and unknown outcomes; the unknown screenshot was an intermediate capture, while its persisted final result was \texttt{REVIEW}. Those historical runs are described separately from the post-merge runs above.

\section*{Limitations and Remaining Work}
The SABLE analyzer supports a bounded subset of Terraform/HCL and IAM/S3 semantics. It does not apply Terraform, query deployed AWS state, or establish the full authorization behavior of AWS. The scenario fixtures are synthetic and authored in-house, so their results are not an independent estimate of real-world effectiveness. SATRA's \texttt{INCONCLUSIVE} status in the retained integrated run also remains visible evidence of an unresolved case, not a passing result.

For FYP-2, the proposed SABLE work is to complete the baseline comparison and broader validation, then extend supported cases only where the research question requires it. The team reports no outstanding proposal or panel actions and confirms that member contributions have been cross-checked. The report records those statements as team confirmations; no signed panel action register, meeting minutes, or approved FYP-2 schedule is represented as documentary evidence. The individual signature lines remain unsigned.
